% Internal research development. The general one-descent theorem is not proved here.
\section{One-sided coefficient comparisons for the reverse kernel}
Write $[t]=1+z+\cdots+z^{t-1}$ and let all sequences be extended by zero
outside their support. For a nonnegative sequence $p$ supported on an
interval beginning at zero, define $\mathcal P_K(p)$ by
\[
 p_{i+1}p_j\ge p_i p_{j+1}\qquad(0\le i\le K,\ j\ge i).
\]
This is a comparison of individual coefficients, not cumulative sums.

\begin{lemma}[Causal convolution preserves a prefix comparison]
If $\mathcal P_K(p)$ holds and $h$ is a finite nonnegative log-concave
sequence with interval support in the nonnegative integers, then
$\mathcal P_K(h*p)$ holds.
\end{lemma}
\begin{proof}
Use upper triangular Toeplitz matrices $T(p)_{ij}=p_{j-i}$. For the
minor of $T(h)T(p)$ in rows $0,1$ and columns $r+1,s+1$, with
$0\le r\le K$ and $s\ge r$, Cauchy--Binet expresses its determinant as
the sum over integers $i<j$ of
\[
 (h_i h_{j-1}-h_j h_{i-1})
 \bigl(p_{r+1-i}p_{s+1-j}-p_{s+1-i}p_{r+1-j}\bigr).
\]
The first factor is nonnegative by log-concavity. A nonzero first
factor requires $i\ge0$. Put $k=r+1-j$, $\ell=s+1-j$, and $d=j-i$.
The second factor is $p_{k+d}p_\ell-p_{\ell+d}p_k$. If $k<0$, this is
nonnegative by zero extension. Otherwise its nonnegativity follows by
multiplying the $d$ successive ratio comparisons at indices
$k,k+1,\ldots,k+d-1$ with the corresponding indices beginning at
$\ell$. All the first indices are at most $r-i\le K$.
Interval support handles zero denominators; equivalently one can
cross-multiply on the positive support and treat endpoints separately.
The sum is finite at the two selected columns. This proves the claim.
\end{proof}


\begin{theorem}[Prefix-to-tail coefficient ratios]
Let $Y$ be weakly increasing, with all parts less than $M\ge2$, let
$b=\operatorname{wt}Y$, and put
\[
 D_Y=1+(S_M-1)F_Y,\qquad \delta_Y(z)=z^{b+M}D_Y(1/z).
\]
Then $\mathcal P_{M-2}(\delta_Y)$ holds. Explicitly, for every integer
$0\le r\le M-2$ and every $s\ge r$,
\[
 \delta_Y(r+1)\delta_Y(s)\ge
 \delta_Y(r)\delta_Y(s+1).
\]
\end{theorem}
\begin{proof}
Induct on the number of parts. The empty arm gives $\delta=[M]$.
For $Y=Zt$, with $c=\operatorname{wt}Z$ and $b=c+t$, use
\[
 \delta_Y=[t]\delta_Z+z^b[M-t].
\]
Call the convolution term $W$. It has $\mathcal P_{M-2}$ by the lemma.
The existing low-rank growth inequality for $D_Y$ says that
$\delta_Y(k+1)\le\delta_Y(k)$ for $k\ge b$.

If $b<M$, then $\delta_Y(k)=A_Y(k):=\sum_{j=0}^kq_Y(j)$ for
$0\le k<M$, where $q_Y=z^bF_Y(1/z)$. The sequence $A_Y$ is
log-concave and is constant from $b$ onward. Thus the required ratios
are ordered inside the prefix. Every later ratio is at most one by
the just-stated monotonicity, whereas the prefix ratios are at least
one. This proves the assertion in this case, including the constant
part of the prefix.

Suppose $b\ge M$. We first show that monotonicity also holds at the
joining edge: $\delta_Y(b)\le\delta_Y(b-1)$. Indeed,
\[
 \delta_Y(b)-\delta_Y(b-1)=1+D_Z(M-t)-D_Z(M).
\]
The difference $D_Z(M)-D_Z(M-t)$ is a nonnegative integer. It is
strictly positive when $b\ge M$: the coefficients of
$(1-z)D_Z$ through rank $M$ equal those of
$E_Z=1-z+(2z-1)F_Z$, whose positive support is
$[\kappa(Z),c+1]$, and $\kappa(Z)\le t+1$. The interval
$[M-t+1,M]$ meets this support, since $c+1=b-t+1\ge M-t+1$ and
$\kappa(Z)\le t+1\le M$. Therefore the displayed difference is at
most zero. This also explains precisely where integrality is used.

For $r\le M-2$ and $s\le b-2$, all four coefficients coincide with
those of $W$, so its prefix comparison applies. For $s\ge b-1$,
the right-hand ratio is at most one, whereas the left-hand ratio is
at least one because $\delta_Y(k)=A_Y(k)$ for $k<M$. Zero endpoints
satisfy the cross-multiplied inequality directly. This completes
the induction.
\end{proof}

\begin{corollary}[A universality-sensitive separation bound]
Let $\kappa=\kappa(Y)$, with $\kappa(\varnothing)=1$. Then
\[
 \delta_Y(k+1)\delta_Y(k+h)\ge
 \delta_Y(k)\delta_Y(k+h+1)
\]
for every integer $k$ and every integer $h\ge b-\kappa+1$.
\end{corollary}
\begin{proof}
The positive support of $\delta_Y$ is $[0,b+M-\kappa]$.
If $k<0$ or $k+h$ is at or beyond the final positive coefficient,
the inequality follows from zero extension. Otherwise
$k+h\le b+M-\kappa-1$, so $k\le M-2$, and the prefix-to-tail
theorem applies. The separation bound is nonnegative because
$\kappa\le b+1$.
\end{proof}

\begin{lemma}[A window preserves log-concavity on the rising side]
Let $p$ be nonnegative, have interval support, be unimodal, and be
log-concave at every index up to its last mode. Then $[t]*p$ has
the same properties for every $t\ge1$. In fact, if $m$ is the last
mode of $p$, the convolution is log-concave at every center
$n\le m+t-1$.
\end{lemma}
\begin{proof}
Unimodality is preserved by a uniform window. Write $q=[t]*p$ and
let $m$ be the last mode of $p$. A comparison centered at $n<m$
uses only the log-concave prefix of $p$, so log-concave convolution
applies after truncating $p$ at $m$. If $n\ge m$ and $n-t<m$,
\[
 (q_n-q_{n-1})-(q_{n+1}-q_n)
   =(p_n-p_{n+1})+(p_{n-t+1}-p_{n-t})\ge0.
\]
Thus $q$ is nonnegative and concave in this region, which implies
log-concavity regardless of the sign of its first difference.
Together with the first case this covers all $n\le m+t-1$.
If $n-t\ge m$, the window
is already nonincreasing; comparisons before its last mode can
only lie on a constant plateau and are immediate. At a mode
the log-concavity inequality is automatic.
\end{proof}

\begin{theorem}[Coefficient log-concavity on the rising side]
For every increasing arm $Y$, both $q_Y=z^bF_Y(1/z)$ and
$\delta_Y$ are unimodal and log-concave at every index up to their
last mode. Equivalently, $F_Y$ and $D_Y$ are log-concave on their
decreasing sides.
\end{theorem}
\begin{proof}
For $q_Y$, use $q_{Zt}=[t]q_Z+z^{c+t}$. The convolution ends at
rank $c+t-1$ with coefficient one. Appending a final coefficient
one cannot affect the rising side, except by extending a terminal
plateau of ones. The window lemma and induction apply.

For $\delta_Y$, the case $b<M$ follows directly from the cumulative
sequence $A_Y$ through $M-1$, followed by a nonincreasing sequence;
its mode plateau includes $[b,M-1]$. If $b\ge M$, use
$\delta_Y=W+z^b[M-t]$ as above. The joining inequality gives
$W_b-W_{b-1}\le-1$, so the last mode of $W$ is at most $b-1$.
The corrected sequence agrees with $W$ before $b$ and is
nonincreasing from $b-1$ onward. The correction therefore cannot
alter a rising-side log-concavity comparison, except possibly one
at a mode itself, which is automatic. The window lemma completes
the induction.
\end{proof}

\begin{corollary}[A coefficient window beyond the initial prefix]
For every $t\ge1$, the coefficient sequence of $[t]\delta_Y$ is
log-concave at every center $n\le M+t-2$.
\end{corollary}
\begin{proof}
The last mode of $\delta_Y$ is at least $M-1$: its first $M$
coefficients are the nondecreasing cumulative coefficients of
$q_Y$. Apply the strengthened window lemma.
\end{proof}

\begin{corollary}[An unconditional terminal likelihood-ratio comparison]
Let $U$ be an increasing arm with all parts at least two, and let
$Y$ be an increasing strict arm for peak $M$. Write $u=\operatorname{wt}U$,
$b=\operatorname{wt}Y$, $T_U=\prod_{t\in U}[t]$, and
$K_U^*=z^u K_U(1/z)$, with the established pure-arm notation. For any
$\alpha\ge0$, set
\[
 L_r=[z^r]K_U^*\delta_Y,\qquad
 B_r=[z^{r+\alpha+1}]T_U\delta_Y.
\]
Then $L_{r+1}B_r-L_rB_{r+1}\ge0$ for every integer $r\le M-2$.
\end{corollary}
\begin{proof}
Put $a_i=[z^i]K_U^*$ and $b_i=[z^{i+\alpha+1}]T_U$.
The pure-arm likelihood-ratio theorem gives
$a_jb_i-a_ib_j\ge0$ for $i<j$. Cauchy--Binet expands the desired
determinant as
\[
 \sum_{i<j}(a_jb_i-a_ib_j)
 \bigl(\delta_{r+1-j}\delta_{r-i}
       -\delta_{r-j}\delta_{r+1-i}\bigr).
\]
A nonzero first factor requires $j\ge0$, since $a$ has nonnegative
support. The lower index $r-j$ of the kernel minor is therefore at
most $M-2$. If it is nonnegative, the prefix-to-tail theorem applies;
if it is negative, zero extension makes the minor nonnegative
directly. All summands are nonnegative.
\end{proof}

\begin{corollary}[A universality-sensitive full comparison]
With the preceding notation, put $s=\operatorname{len}U$.
If $\alpha\ge b+u-s-\kappa(Y)$, then the likelihood-ratio
determinants are nonnegative for every integer $r$, and
$F_{1^\alpha U}D_Y$ is unimodal.
\end{corollary}
\begin{proof}
In a nonzero Cauchy--Binet term, $j\ge0$ and
$i\le u-s-\alpha-1$, hence
$j-i\ge\alpha+1-(u-s)\ge b-\kappa(Y)+1$.
The separation corollary makes every kernel minor nonnegative.
Interval support of the denominator then gives the global
single-crossing property in the exact difference identity for
$F_{1^\alpha U}D_Y$.
\end{proof}

\begin{theorem}[A light non-one skeleton]
Let $X=1^\alpha U$ and $Y$ be increasing strict arms for peak $M\ge2$,
where $U$ is nonempty and has parts at least two. Put
$a=\operatorname{wt}X$, $b=\operatorname{wt}Y$, and
$u=\operatorname{wt}U$. If
\[
 M+a-b\ge2u-2,
\]
then both $F_XD_Y$ and $F_{XMY}$ are unimodal. In particular a heavy arm with
$u\le(M+2)/2$ is covered with no restriction on arm lengths.
\end{theorem}
\begin{proof}
Put $N=a+b+M$. The existing interval-layer reflection for
$Q=F_XD_Y$ gives $\Delta Q(n)\ge0$ for
$n\le\lfloor(N+2)/2\rfloor$. The hypothesis and $u\ge2$ imply
$N\ge2b+2$, so the universal first-half growth of $H=F_Y+Q$
covers $n\le b+1$.
For $n\ge b+2$ the term $F_Y$ contributes no difference. Thus
$\Delta H(n)\ge0$ throughout $n\le\lfloor(N+2)/2\rfloor$.

The bridge identity gives
$\Delta H(n)=L_r-B_r$ for $r=u+b+M-n$ and $n>b+1$.
The preceding corollary orders the ratios $L_r/B_r$ on all indices
$r\le M-1$. The denominator has interval support; before that
support $L=0$, and beyond it $L-B\ge0$. Consequently these
differences cannot change from positive to negative as $r$ increases
within that interval. Reversing the index shows that $\Delta H(n)$
cannot change from negative to positive for $n\ge u+b+1$.
This entire range satisfies $n>b+1$, since $u\ge2$.

The parameter inequality is equivalent to $N\ge2u+2b-2$, hence
$\lfloor(N+2)/2\rfloor\ge u+b$. The increasing portion therefore
joins the terminal single-crossing portion without a missing
difference. This proves unimodality.
The same proof applies directly to $Q$, whose exact difference
identity has no $F_Y$ correction at any rank.
\end{proof}

\begin{corollary}[Arbitrary peaks for a two or a three after leading ones]
Every strict core with one arm of the form $1^\alpha 2$ or
$1^\alpha 3$ has a unimodal rank polynomial, with no bound on
$M$, $\alpha$, or the other arm. The corresponding relative
polynomials are unimodal as well.
\end{corollary}
\begin{proof}
Fix $M$ and consider the family of ordered arm pairs for which at
least one arm has total non-one weight at most three. This family
is closed under removing a last arm part and under exchanging
the arms. Therefore the balanced-minimal-counterexample proof
and the outer-three-layer proof apply within this family: every
shorter pair required in their inductions remains in the family.
The empty and rectangular-arm bases are unchanged. A minimal
counterexample would satisfy $|a-b|\le M-4$ for $M\ge3$;
the case $M=2$ is already a rectangular-arm case. Orient the pair
so that $X$ has non-one weight $u\le3$. The rectangular case
$u=0$ is excluded. Thus $2\le u\le3$, and
$M+a-b\ge4\ge2u-2$. The preceding theorem excludes the
putative counterexample, so all relative polynomials in the
family are unimodal. For a core, orient the heavier arm on the
left and use the established core/relative equivalence; the
family is unchanged by this orientation. This proves the claim.
\end{proof}

\paragraph{Exact remaining limitation.}
These results compare actual coefficients and strengthen the cumulative
inequalities. They do not prove log-concavity of $F_XD_Y$, nor do they
prove the restricted heavy-arm likelihood-ratio assertion through
$r\le\operatorname{wt}(U)+M-3$. Products of unimodal sequences whose
rising sides are log-concave need not be unimodal. For instance,
$(1,3,2,1,1,1,1)$ has these properties, whereas its self-convolution has
coefficients $(1,6,13,14,12,12,13,12,7,4,3,2,1)$. The composition-specific
product step remains to be established.

\paragraph{A mode-sensitive version of the terminal comparison.}
The rising-side theorem yields a larger, data-dependent prefix.
For any nonnegative interval-supported sequence which is unimodal
and log-concave up to its last mode $m$, the property
$\mathcal P_{m-1}$ holds: comparisons before $m$ follow by
successive log-concavity, and comparisons with later indices compare
a ratio at least one with a ratio at most one. Consequently the
terminal likelihood-ratio proof applies to every $r\le m_Y-1$,
where $m_Y$ is the last mode of $\delta_Y$.

Let $\ell=\operatorname{len}Y$. Its positive expansion is
\[
 \delta_Y=[M]\prod_{j=1}^{\ell}[y_j]
  +\sum_{j=1}^{\ell}z^{y_1+\cdots+y_j}[M-y_j]
       \prod_{i=j+1}^{\ell}[y_i].
\]
Put $d=b-\ell+M-1$. The first summand is symmetric about $d/2$,
and the $j$th additional summand has center
$(d+y_1+\cdots+y_{j-1}+j)/2\ge d/2$.
Hence
\[
 m_Y\ge \max\left(M-1,\left\lceil
                  \frac{b-\ell+M-1}{2}\right\rceil\right).
\]
The exact difference identity and the mode-sensitive comparison
show that $Q_{1^\alpha U,Y}$ is unimodal whenever
\[
 \left\lfloor\frac{a+b+M+2}{2}\right\rfloor+m_Y
       \ge u+b+M-1.
\]
Indeed the first-half growth covers every rank before
$u+b+M-m_Y$, and the likelihood-ratio argument rules out a
negative-to-positive change thereafter.

\begin{theorem}[One long initial block of ones]
Suppose one of the two increasing strict arms is $X=1^\alpha U$,
where every part of $U$ is at least two, and put
$u=\operatorname{wt}U$, $\ell=\operatorname{len}Y$. If
\[
 \alpha\ge u+\ell-3,
\]
then the two-arm rank polynomial $F_{XMY}$ is unimodal.
Both relative polynomials $F_XD_M(Y)$ and $F_YD_M(X)$ are also
unimodal. No lower bound on the leading ones in the other arm,
or on the relative arm weights, is required. An all-one arm is
included independently of the displayed inequality.
\end{theorem}
\begin{proof}
First orient the qualifying arm as $X$. If $U$ is empty, its
relative polynomial is a uniform window of the unimodal $D_M(Y)$.
Otherwise put
$g=\lfloor(a+b+M+2)/2\rfloor$ and
$k=\lceil(b-\ell+M-1)/2\rceil$. Since $a=\alpha+u$,
\[
 g+k\ge b+M+\left\lfloor\frac{a-\ell+1}{2}\right\rfloor
      \ge b+M+u-1.
\]
The mode-sensitive criterion proves unimodality of the relative
polynomial in this orientation.

To obtain both orientations and the rank polynomial, fix $M$ and
consider the family of pairs in which at least one arm satisfies
the stated condition or is all ones. This family is closed under
swapping arms and deleting a final part from either arm. Deleting
a part from the qualifying non-one skeleton decreases $u$;
deleting a part from the other arm decreases $\ell$. If a skeleton
becomes empty, the all-one case applies.

If a relative counterexample existed in this family, choose one
with least total number of arm parts. The established least-relative-
counterexample argument, which uses only these deletions, gives
$|a-b|\le M-1$. In this band the rank polynomial and each of the
two relative polynomials have equivalent unimodality. Orienting
the qualifying arm as $X$ contradicts the result just proved.
Thus both relative orientations are unimodal throughout the
family. Finally orient a heavier arm as the first arm; the
rank/relative equivalence applies because $b\le a+M-1$.
It gives the asserted rank unimodality.
\end{proof}

\begin{proposition}[Grouping the initial ones improves the mode bound]
Write $Y=1^\beta V$, where $V=(t_1,\ldots,t_s)$ has parts at least
two, and set $v=\operatorname{wt}V$, $e=v-s$, $b=\beta+v$.
For every strict peak $M$,
\[
 m_Y\ge \max\left(M-1,
   \left\lceil\frac{e+M-1}{2}\right\rceil,
   \left\lceil\frac{e+M-2+\min(\beta,M)}{2}\right\rceil\right).
\]
\end{proposition}
\begin{proof}
The first two bounds were proved above. The all-one kernel has the
exact identity
\[
 \delta_M(1^\beta)=[M-1][\beta+1]+z^{M-1}.
\]
Writing $T_V=\prod_j[t_j]$ and $v_j=t_1+\cdots+t_j$, repeated
use of the positive kernel recurrence gives
\[
 \delta_M(Y)=T_V[M-1][\beta+1]+z^{M-1}T_V
 +\sum_{j=1}^s z^{\beta+v_j}[M-t_j]
                       \prod_{i>j}[t_i].
\]
These are nonnegative shifted symmetric unimodal interval products.
Twice their centers are, respectively,
\[
 e+M-2+\beta,\qquad e+2M-2,\qquad
 e+M-1+2\beta+v_{j-1}+j.
\]
Every center is at least
$(e+M-2+\min(\beta,M))/2$. Each summand is therefore
nondecreasing through the ceiling of that number. Their sum has
the same property. Since the kernel is already known to be unimodal,
its last mode is at least this ceiling.
\end{proof}

\begin{theorem}[A stronger initial-one criterion]
Let $X=1^\alpha U$ and $Y=1^\beta V$ be the two strict increasing
arms, where the displayed skeletons contain only parts at least two.
Put $u=\operatorname{wt}U$ and $s=\operatorname{len}V$. If
\[
 \alpha\ge u+s+\max(\beta-M,0)-2,
\]
then the rank polynomial and both relative polynomials are unimodal.
The criterion can be applied in either orientation. All-one arms
remain unconditional base cases. For $\beta=0$, the preceding
criterion $\alpha\ge u+s-3$ is slightly stronger and still applies.
\end{theorem}
\begin{proof}
Set $r=\max(\beta-M,0)$, so $\min(\beta,M)=\beta-r$.
The new mode bound is
\[
 m_Y\ge\left\lceil\frac{b-s+M-2-r}{2}\right\rceil.
\]
For $a=\alpha+u$, it yields
\[
 \left\lfloor\frac{a+b+M+2}{2}\right\rfloor+m_Y
 \ge b+M+\left\lfloor\frac{\alpha+u-s-r}{2}\right\rfloor
 \ge u+b+M-1.
\]
Thus the mode-sensitive terminal bridge proves unimodality of
$F_XD_M(Y)$ in this orientation. The family of pairs satisfying
the stated condition in at least one orientation, together with
the all-one-arm pairs, is closed under swapping arms and deleting
a last part from either arm. Deleting a non-one part of $U$
decreases $u$, and deleting a non-one part of $V$ decreases $s$;
deletions of leading ones occur only after that arm becomes all ones.
The same least-relative-counterexample argument used in the
previous theorem therefore gives both orientations and the full
rank polynomial. This proves a sufficient condition; it does not
resolve the unrestricted two-skeleton case.
\end{proof}

\begin{corollary}[One non-one part in each arm, with arbitrary padding]
For every $M>\max(t,v)\ge2$ and all $\alpha,\beta\ge0$, the
core
\[
 (1^\alpha,t,M,1^\beta,v)
\]
has a unimodal rank polynomial. Both relative polynomials are
unimodal as well. The peak and both padding lengths are unbounded.
\end{corollary}
\begin{proof}
An arm with zero padding consists of one part and has the same
rank polynomial as an all-one arm of the same weight. Such a case
is already covered. For positive padding, the exact identity
\[
 F_{1^\alpha t}=1+z[t][\alpha+1]
\]
permits exchange of the two interval factors. Consequently each
arm can be replaced, without changing its weight or rank polynomial,
by $1^{\alpha'}t'$ with $2\le t'<M$ and $\alpha'\ge t'-1$.
The replacement is needed only when $\alpha<t-1$, in which case
$(\alpha',t')=(t-1,\alpha+1)$. Normalize both arms in this way.

Choose the orientation with $\alpha'\ge\beta'$. If $\beta'\le M$,
then $\alpha'\ge t'-1$. If $\beta'>M$, then
$\alpha'-(\beta'-M)\ge M>t'-1$. In either case
\[
 \alpha'\ge t'+1+\max(\beta'-M,0)-2.
\]
This is the stronger initial-one criterion with a one-part opposite
skeleton. Its conclusion transfers back to the original pair,
because the core and both relative polynomials depend only on the
two arm polynomials. The zero-padding cases are transferred in
the same way from the all-one-arm base theorem.
\end{proof}

\begin{corollary}[A run tied exactly to the peak]
For $2\le t<M$, a core with one arm $1^{M-1}t$ is unimodal
for every increasing opposite arm. Both relative polynomials
are unimodal as well.
\end{corollary}
\begin{proof}
The initial exchange gives
$F_{1^{M-1}t}=1+z[t][M]=F_{(t-1),M}$. Substitution into the
maximum decomposition gives the exact identity
\[
 F_{1^{M-1}t,M,Y}=1+S_M F_{(t-1),M,Y}.
\]
The shorter core has a single-part arm, so its rank polynomial
is unimodal by the established rectangular-arm theorem.
Uniform convolution followed by adding the initial constant
preserves unimodality. For both relative orientations, apply the
least-counterexample argument within the family having such an
arm, including the all-one-arm base cases. This family is closed
under arm exchange and final-part deletions, and every full core
in it has just been proved unimodal. A least relative counterexample
would be balanced and hence would give a nonunimodal full core,
a contradiction.
\end{proof}

